\chapter*{Resumen}
\addcontentsline{toc}{chapter}{Resumen}

El presente anteproyecto de graduación, titulado \textit{Diseño de un sistema de bajo costo 
para el control de ESD para laboratorios de electrónica}, propone el diseño e implementación 
de una estación orientada a mejorar la trazabilidad de las pruebas de descarga 
electrostática (ESD) en entornos de laboratorio. El proyecto surge a partir de una oportunidad 
de mejora identificada en el laboratorio TEM de Intel Costa Rica, donde actualmente se 
utiliza un probador comercial de pulseras antiestáticas, pero no se dispone de un sistema 
integrado que asocie cada prueba con un usuario, almacene registros históricos y genere 
reportes de cumplimiento.

La solución propuesta integra un probador comercial de pulseras antiestáticas como subsistema 
principal de verificación ESD, identificación de usuarios mediante lectura de código de barras 
y reconocimiento facial, interfaz gráfica táctil, monitoreo de temperatura y humedad relativa, 
almacenamiento local de registros y generación de reportes históricos. Además, se plantea una 
arquitectura modular y personalizable con un costo estimado aproximado de 250 USD, valor menor 
al de soluciones comerciales avanzadas analizadas en el estado del arte. La validación del 
prototipo se realizará mediante pruebas funcionales de los subsistemas, pruebas de integración 
y pruebas de operación en el laboratorio TEM de Intel Costa Rica. El proyecto no pretende 
sustituir sistemas industriales certificados, sino demostrar la viabilidad de una estación ESD 
, accesible y adaptable para mejorar el seguimiento del cumplimiento de procedimientos 
ESD en laboratorios de electrónica.

\textbf{Palabras clave:} ESD, descarga electrostática, control ESD, pulsera antiestática, 
trazabilidad, laboratorio de electrónica, semiconductores, Intel TEM, bajo costo, 
interfaz gráfica, identificación de usuario, reconocimiento facial, código de barras, 
monitoreo ambiental, registros históricos.


\chapter*{Abstract}
\addcontentsline{toc}{chapter}{Abstract}

This graduation project proposal, entitled \textit{Design of a Low-Cost ESD Control System 
for Electronics Laboratories}, proposes the design and implementation of an intelligent 
station intended to improve the traceability of electrostatic discharge (ESD) tests in 
laboratory environments. The project arises from an improvement opportunity identified 
at the Intel Costa Rica TEM Laboratory, where a commercial wrist strap tester is 
currently used, but no integrated system is available to associate each test with a user, 
store historical records, and generate compliance reports.

The proposed solution integrates a commercial wrist strap tester as the main ESD verification 
subsystem, user identification through barcode reading and facial recognition, a touchscreen 
graphical interface, temperature and relative humidity monitoring, local data logging, and 
historical report generation. In addition, a modular and customizable architecture is proposed 
with an estimated cost of approximately USD 250, which is lower than the advanced commercial 
solutions analyzed in the state of the art. The prototype validation will be performed through 
functional subsystem tests, integration tests, and operational tests at the Intel Costa Rica 
TEM Laboratory. The project does not aim to replace certified industrial ESD systems, but 
rather to demonstrate the feasibility of an intelligent, accessible, and adaptable ESD station 
for improving ESD procedure compliance tracking in electronics laboratories.

\textbf{Keywords:} ESD, electrostatic discharge, ESD control, wrist strap, traceability, 
electronics laboratory, semiconductors, Intel TEM, low cost, graphical interface, 
user identification, facial recognition, barcode, environmental monitoring, historical records.
